\vspace{-0.1cm}
\section{Conclusion}
\vspace{-0.1cm}

In this paper we have presented a novel segmentation algorithm for wind turbine blade images. To this end, we have proposed a BU-Net model in combination with additional steps based on hole filling and random forest strategies that can solve the problem in an accurate and efficient manner. Our BU-Net exploits a focal loss that is regularized with a contiguity term. The hole filling is simple but effective, and helps to improve the solution after applying every learning model. Finally, our unsupervised random forest is trained on-the-fly capturing the main patterns, like a denoising step. Remarkable quantitative and qualitative results are provided on newly acquired blade images, even when they belong to distinct inspection campaigns and wind farms, validating the capability of our model to generalize properly. Our future work is oriented to extend our model to semantic segmentation where different parts of the wind turbine need to be segmented in the image. Furthermore, we would like to study its impact on facilitating the defect detection task.

%Thanks to our segmentation algorithm, we believe our framework has the potential to be used for many applications for structural health monitoring. From a wind industry perspective, that means we can start to automate blade maintenance, saving costs and making green energy more competitive compared to traditional fuel power.

%This work exhibits the different strategies that were carried out to accomplish a novel image segmentation algorithm for wind turbine blade images. That is, a customized U-Net model complemented by two main post-processing approaches: an extended hole filling procedure and an unsupervised random forest trained on the fly. In particular, we regularized the focal loss with a contiguity term and we employed rotation-based data augmentation to define our BU-Net deep-learning model.

%Remarkable results are produced on new acquired blade images, even when they belong to distinct inspection campaigns and windfarms. This proves its capability to generalize properly. Furthermore, we believe our framework has the potential to be used for many applications for structural health monitoring, starting with facilitating the defect detection task. From a wind industry perspective, that means we can start to automate blade maintenance, saving costs and making green energy more competitive compared to traditional fuel power.






%compared to traditional fuel power
%submitted
%variants 
%currently the 
%top performing method on the PASCAL
%report
%conducted 
%attaining
%focus
%yields
%significantly outperforming
%further improves performance to 63.1 percent. 
%observe any improvement when adopting either
%Employing the dense
%substantially outperforms the concurrent work 
%examine
%cover
%introduce 
%measure
%suggest
%To this aim
%OPERATIONS MAINTENANCE